\subsection{Research object and forecasting task}\label{sec:forecast-task}
The empirical objects are timestamped geographic position recordings from
a historical GPX collection. Each recording traces the movement of an
unidentified recording carrier. The retained archive does not establish
whether a carrier is walking, cycling or travelling by another means, or
whether two recordings belong to the same person. We therefore study
position changes in the selected recordings, rather than the behaviour of
a defined participant population. A cleaned segment is a derived portion
of a recording; a window splits that portion into an observed prefix and
future targets. Multiple windows can belong to one recording-hash block.
Only one primary origin per selected block enters the main evaluation.

For an origin with saved time $t_0$, let $\mathcal H_0$ contain the allowed
observed position/time prefix, ending at the current position $x_0$.
Its last three pairs initialize the retained motion-history buffer.
The task is to predict future positions conditional on this visible
information and on fixed development-fitted parameters $\widehat\xi$.
The method's internal east/north frame is anchored at the first visible
position of that whole prefix, which need not be the first of the last
three buffered points. Terrain prediction and evaluation instead use a frame anchored
at the last visible position, the forecast origin. Method outputs are
converted through geographic coordinates into that common scoring frame;
the distributions in Equation~\eqref{eq:forecast-object} are expressed
there. The method also has the retained clock and
its solar proxy; the terrain branch additionally queries fixed map fields
at simulated positions. True future positions are reserved for selection
and evaluation and do not update a forecast. With actual matched elapsed
times $u_1,\ldots,u_4$, the reported predictive object is
\begin{equation}\label{eq:forecast-object}
 \widehat P_j=\mathcal L_{\widehat\xi}(X_{t_0+u_j}\mid\mathcal H_0),
 \qquad
 \widehat P_{N,j}=\frac1N\sum_{p=1}^{N}\delta_{X_j^{(p)}},
 \quad j=1,\ldots,4 .
\end{equation}
Here $\delta_x$ is a point mass at $x$, and each simulated particle $p$
has a position at all four targets. The ensemble approximates a
model-defined distribution; it is not a sample of people. Its mean gives
a representative location, while its radial quantiles define prediction
disks. These summaries answer different questions about the same forecast.
No separate measurement-error model is fitted: the saved target position
is used as the evaluation reference, although its physical accuracy is
not independently verified.

Development windows supply parameter estimates; final-evaluation windows
supply forecast origins and observed outcomes. The main analysis asks how
well those fixed predictors describe these outcomes as the saved forecast
horizon increases. It does not estimate a causal law of human movement or
the probability of finding a person in an operational search.
